\documentclass[11pt]{article}
\usepackage[T1]{fontenc}
\usepackage{lmodern,amsmath,amssymb,amsthm,mathtools,booktabs}
\usepackage[margin=1in]{geometry}
\usepackage{microtype}
\usepackage[colorlinks=true,linkcolor=blue,citecolor=blue,urlcolor=blue]{hyperref}
\usepackage{enumitem}
\setlist{itemsep=2pt,topsep=4pt}
\newtheorem{theorem}{Theorem}[section]
\newtheorem{lemma}[theorem]{Lemma}
\theoremstyle{remark}
\newtheorem{remark}[theorem]{Remark}
\numberwithin{equation}{section}
\title{Orbit Constraints in Symmetric Arrays\\and Circulant Weighing Matrices}
\author{}
\date{Research draft -- 24 October 2026}
\hypersetup{pdftitle={Orbit Constraints in Symmetric Arrays and Circulant Weighing Matrices},pdfsubject={Catalan-bound tightness and nonexistence of CW(195,81)}}
\begin{document}
\maketitle
\begin{abstract}
We establish two results in combinatorial design theory. For every $k\ge3$,
the minimum number of rows in a binary strength-two orthogonal array invariant
under all coordinate permutations is
$4\binom{k-2}{\lfloor(k-2)/2\rfloor}$, proving Karunanayaka's
Catalan-bound tightness conjecture. We also classify all minimizing shell
multiplicity vectors. Separately, we prove that a circulant weighing matrix
$CW(195,81)$ does not exist, resolving a parameter pair listed as open in
Arasu--Gordon--Zhang's table. Both proofs are analytic. The first combines
indivisible shell masses with two moment constraints. The second combines
contraction, established cyclotomic multiplier theory, a kernel character,
and a short augmentation obstruction. Exact finite checks are supplied as
independent corroboration, not as premises of either proof.
\end{abstract}

\section{Scope and relation to the proposed method}

The purpose of this note is to test a proof strategy outside periodic
Dirichlet-series questions: compress constraints by symmetry, retain the
integrality and coefficient restrictions that survive compression, and use
them to force a sharp bound or a contradiction. The results below concern
different mathematical objects and do not assume any claim about an
Erd\H{o}s function or an Okada lattice.

The first target is Conjecture~1 of Karunanayaka~\cite{Karunanayaka2026},
whose published statement covers every $k\ge3$ and whose reported finite
verification covers $k\le14$. We prove the original statement with its full
coordinate-permutation symmetry assumption. The second target is the
existence of $CW(195,81)$, listed in Table~10 of
Arasu--Gordon--Zhang~\cite{ArasuGordonZhang2021}.

The tools used here have established precedents. Cauchy--Schwarz and orbit
compression are standard. The cyclotomic multiplier theorem is due to
Tan~\cite{Tan2018}; kernel-character reductions also appear in
Tan~\cite{Tan2026}. The contributions asserted in this draft are the
specific lower-bound proof and equality classification, and the elimination
of the parameter pair $(195,81)$. No new general multiplier theorem or
universal $2$-adic decision procedure is asserted.

\paragraph{Review status.}
The proofs have undergone separate internal mathematical audits and exact
arithmetic checks. This is an AI-assisted research draft, not an externally
refereed publication. The source check on 24 October 2026 found the targets
stated as open in the cited records and found no later resolution of these
specific statements. That search is not a guarantee of publication priority.

\section{The Catalan bound for fully symmetric binary arrays}
\label{sec:catalan-bound}

A binary array is \emph{fully symmetric} if its multiset of rows is invariant
under every permutation of the coordinates.  Equivalently, there are
nonnegative integers $m_0,\ldots,m_k$ such that every binary word of weight
$w$ occurs exactly $m_w$ times.  Repeated rows are allowed.
Karunanayaka~\cite{Karunanayaka2026} constructs fully symmetric strength-two
orthogonal arrays with
\[
 B_k=4\binom{k-2}{\lfloor(k-2)/2\rfloor}
\]
rows and conjectures that this is the minimum for every $k\ge3$.
We prove this lower bound and determine all equality cases.

\begin{theorem}\label{thm:catalan-minimum}
Every nonempty fully symmetric binary orthogonal array of strength two on
$k\ge3$ coordinates has at least $B_k$ rows.  Equality holds precisely for
the following multiplicity vectors, with all unlisted entries zero:
\begin{enumerate}
\item If $k=2r+1$ is odd, either $m_{r+1}=1$, $m_0=C_r$, or its level flip
$m_r=1$, $m_k=C_r$.
\item If $k=2r\ge6$ is even, $m_r=1$ and $m_0=m_k=C_{r-1}$.
\item If $k=4$, the two vectors are $(1,0,1,0,1)$ and $(0,1,0,1,0)$.
\end{enumerate}
Here $C_r=\binom{2r}{r}/(r+1)$ is the $r$th Catalan number.
\end{theorem}

The proof uses two moments and the fact that a populated shell contains at
least one full coordinate-permutation orbit.

\begin{lemma}[A forced-shell bound]\label{lem:forced-shell}
Let a multiset of $N>0$ real numbers have sum zero and square-sum $kN$, with
$k>0$.  Suppose $M$ of its entries are equal to $d$.  Then
\[
 N\ge M\left(1+\frac{d^2}{k}\right).
\]
If also $d=0$ and every entry has absolute value at most $k>1$, then
$N\ge kM/(k-1)$.
\end{lemma}

\begin{proof}
The remaining $N-M$ entries have sum $-Md$ and square-sum $kN-Md^2$.
Cauchy--Schwarz gives
\[
 M^2d^2\le(N-M)(kN-Md^2),
\]
which simplifies to the first inequality.
For the second, $kN\le k^2(N-M)$.
\end{proof}

\begin{lemma}\label{lem:binomial-shell}
For $k\ge3$ and nonzero integers $d$ with $d\equiv k\pmod2$ and $d^2\le k$,
\[
 \binom{k}{(k-|d|)/2}\left(1+\frac{d^2}{k}\right)\ge B_k.
\]
If $k$ is odd, equality holds only for $|d|=1$.
If $k$ is even, equality holds only when $k=4$ and $|d|=2$.
\end{lemma}

\begin{proof}
For nonnegative admissible $d$ put
$G_k(d)=\binom{k}{(k-d)/2}(k+d^2)$.
The ratio of consecutive values is
\[
 \frac{G_k(d+2)}{G_k(d)}
 =\frac{k-d}{k+d+2}\frac{k+(d+2)^2}{k+d^2}.
\]
The difference of cross products is
\[
 (k-d)(k+(d+2)^2)-(k+d+2)(k+d^2)
 =2(d+1)\bigl(k+1-(d+1)^2\bigr).
\]
This is positive whenever $(d+2)^2\le k$, so $G_k$ strictly increases along
the admissible range.
For odd $k$, its least nonzero argument is $1$, and
\[
 \frac{G_k(1)}k
 =\frac{k+1}{k}\binom{k}{(k-1)/2}=B_k.
\]
For even $k$, the least nonzero argument is $2$.  With
$C=\binom{k}{k/2}$, we have
\[
 \frac{G_k(2)}k=C\frac{k+4}{k+2}
 \ge C\frac{k}{k-1}=B_k,
\]
where the difference of the two rational factors is
$(k-4)/((k+2)(k-1))$.
This also gives the equality statements.
\end{proof}

\begin{proof}[Proof of Theorem~\ref{thm:catalan-minimum}]
Replace each binary coordinate $x_i$ by $1-2x_i$, and let $D$ denote the sum
of the $k$ signs in a row.  Strength two gives
\[
 \sum D=0,\qquad \sum D^2=kN.
\]
In the shell of weight $w$, the value of $D$ is $d_w=k-2w$ and the number of
rows is $M_w=m_w\binom{k}{w}$.
Since the average of $D^2$ is $k$, some populated shell has $d_w^2\le k$.
Write $d=d_w$ and $M=M_w$ for such a shell.  Its positive integer
multiplicity gives
\[
 M\ge\binom{k}{(k-|d|)/2}.
\]
If $d\ne0$, Lemmas~\ref{lem:forced-shell} and~\ref{lem:binomial-shell}
imply $N\ge B_k$.
If $d=0$, then $k$ is even and Lemma~\ref{lem:forced-shell} gives
\[
 N\ge\frac{k}{k-1}M
 \ge\frac{k}{k-1}\binom{k}{k/2}=B_k.
\]
This covers the boundary $d^2=k$ as well.

For equality at odd $k$, Lemma~\ref{lem:binomial-shell} forces $|d|=1$,
and the selected shell has multiplicity one.
Equality in Cauchy--Schwarz makes all remaining values of $D$ equal.
Their mean is $-k/d$, so they are all the opposite endpoint $\mp k$.
The endpoint multiplicity is
\[
 \frac1k\binom{k}{(k-1)/2}=C_{(k-1)/2},
\]
giving the two displayed vectors.

For even $k\ge6$, the nonzero-$d$ bound is strict, so equality requires a
central shell of multiplicity one.
Equality in $\sum D^2\le k^2(N-M)$ makes all remaining rows endpoints.
Zero mean gives equal endpoint multiplicities, each
\[
 \frac1{2(k-1)}\binom{k}{k/2}=C_{k/2-1}.
\]
For $k=4$ this still gives $(1,0,1,0,1)$; equality with $|d|=2$ is also
possible, and Cauchy--Schwarz then forces the opposite shell of weight
$3$ or $1$, giving $(0,1,0,1,0)$.

Finally, every displayed vector satisfies $\sum D=0$ and $\sum D^2=kN$.
Coordinate symmetry turns these two identities into zero first-order sign
means and zero pairwise sign correlations, hence strength two.
Their row counts are $B_k$, proving attainability and the complete
classification.
\end{proof}

\begin{remark}
The upper-bound constructions are those of
Karunanayaka~\cite{Karunanayaka2026}.  The lower bound uses integrality only
through the implication $m_w>0\Rightarrow m_w\ge1$.
This is a theorem for full coordinate-permutation symmetry, not for
arbitrary binary orthogonal arrays.
\end{remark}

\begin{remark}
For $k=15$, the minimum is $B_{15}=6864$.
The two minimizers have $m_8=1$, $m_0=429$, and their level flip
$m_7=1$, $m_{15}=429$; indeed $\binom{15}{8}+429=6435+429=6864$.
For $k=16$, the minimum is $B_{16}=13728$, and the unique minimizer has
$m_8=1$, $m_0=m_{16}=429$, since
$\binom{16}{8}+2\cdot429=12870+858=13728$.
These are the first two values beyond the range $k\le14$ verified in
\cite{Karunanayaka2026}.
\end{remark}

\section{Nonexistence of a circulant weighing matrix of order 195 and weight 81}
\label{sec:cw195}

A circulant weighing matrix $CW(v,k)$ is equivalent to an element
$A\in\mathbb Z[C_v]$ whose coefficients belong to $\{0,1,-1\}$ and which
satisfies $AA^{(-1)}=k$. Here the involution combines inversion of group
elements with complex conjugation of coefficients. The augmentation
$\epsilon(A)$ is the sum of the coefficients. We prove the following
parameter-specific nonexistence result.

\begin{theorem}\label{thm:cw195}
There is no circulant weighing matrix $CW(195,81)$.
\end{theorem}

The pair $(195,81)$ appears among the unresolved cases in
Arasu--Gordon--Zhang, \emph{New nonexistence results on circulant weighing
matrices}, Table~10, arXiv:1908.08447v3 (2021). Their Table~9 records two
equivalence classes of $ICW_3(65,81)$. Tan's September 2026 preprint,
\emph{Character and Multiplier Obstructions for Circulant Weighing Matrices},
arXiv:2608.18468v2, resolves eight other entries. Its methods include the
same general kernel-character reduction employed below; neither that
reduction nor the multiplier theorem is claimed as new here.

We use the integral multiplier theorem and the following cyclotomic version,
proved in Tan, \emph{Group Invariant Weighing Matrices},
arXiv:1610.01914v4 (2018), Theorem~3.4 and Corollary~3.5(b).
If $X\in\mathbb Z[\zeta_u][G]$, where $G$ is finite abelian,
\[
 XX^{(-1)}=n,\qquad (n,|G|)=1,
\]
and the cyclotomic automorphism $\sigma_t$ fixes every prime ideal above
the prime divisors of $n$ in $\mathbb Q(\zeta_{\operatorname{lcm}(u,\exp G)})$,
then $X^{(t)}=\eta gX$ for a root of unity $\eta\in\mathbb Z[\zeta_u]$
and $g\in G$. If additionally $(t-1,\exp G)=1$, a group translation
removes $g$. In $X^{(t)}$, the coefficients are acted on by $\sigma_t$
and the group elements are raised to their $t$th powers.

\begin{lemma}\label{lem:icw65-divisible}
Every $B\in\mathbb Z[C_{65}]$ satisfying
\[
 BB^{(-1)}=81,\qquad \epsilon(B)=9,
 \qquad |B_x|\le3
\]
has coefficients in $\{0,3,-3\}$.
\end{lemma}

\begin{proof}
Apply Tan's generalized multiplier theorem with coefficient conductor $u=1$,
group $C_{65}$, norm $n=81$, and multiplier integer $3$.
Here $(81,65)=1$, and the automorphism $\zeta_{65}\mapsto\zeta_{65}^3$
is the Frobenius at $3$, fixing every prime ideal above $3$.
Thus $B^{(3)}=\eta X^jB$ with $\eta\in\{1,-1\}$.
Augmentation $\epsilon(B)=9$ forces $\eta=1$, and $(3-1,65)=1$
permits a group translation with $B^{(3)}=B$.
This applies to the bounded integer contraction directly; it does not require
its coefficients to be $0,\pm1$.
Use $C_{65}=C_5\times C_{13}$. Multiplication by $3$ has order $4$ on
$C_5\setminus\{0\}$ and order $3$ on $C_{13}\setminus\{0\}$.
Its orbit coefficients can therefore be written as
\[
 a\quad(\text{size }1),\qquad b\quad(\text{size }4),\qquad
 c_1,\ldots,c_4\quad(\text{size }3),\qquad
 d_1,\ldots,d_4\quad(\text{size }12).
\]
The $c_i$ and $d_i$ use the same ordering of the four nonzero $C_{13}$
orbits. All ten coefficients are integers of absolute value at most $3$.
The energy equation is
\begin{equation}\label{eq:cw65-energy}
 a^2+4b^2+3\sum_i c_i^2+12\sum_i d_i^2=81.
\end{equation}
Projecting onto $C_5$ gives coefficient
$x=a+3\sum_i c_i$ at zero and coefficient
$y=b+3\sum_i d_i$ at each nonzero element. Consequently
\[
 x+4y=9,\qquad x^2+4y^2=81.
\]
Elimination gives $4y(5y-18)=0$, whence $y=0$ and $x=9$ by integrality.
Thus $a=3\alpha$ and $b=3\beta$ with
$\alpha,\beta\in\{-1,0,1\}$, and
\begin{equation}\label{eq:cw65-sums}
 \sum_i c_i=3-\alpha,\qquad \sum_i d_i=-\beta.
\end{equation}

The projection onto $C_{13}$ has energy
\begin{equation}\label{eq:cw65-energy13}
 (a+4b)^2+3\sum_i(c_i+4d_i)^2=81.
\end{equation}
If $\beta\ne0$, then $|3\alpha+12\beta|\le9$ forces
$\alpha=-\beta$. The first term of \eqref{eq:cw65-energy13} is then
$81$, so $c_i+4d_i=0$ for all $i$. Since $|c_i|\le3$, every $d_i=0$,
contradicting \eqref{eq:cw65-sums}. Hence $\beta=0$.

Set $t=\sum_i d_i^2$. Comparing
\eqref{eq:cw65-energy} and \eqref{eq:cw65-energy13} now yields
\begin{equation}\label{eq:cw65-reduced}
 \sum_i d_i=0,\qquad \sum_i c_i=3-\alpha,
 \qquad \sum_i c_i^2+4t=27-3\alpha^2,
 \qquad 2\sum_i c_id_i+3t=0.
\end{equation}
For $t>0$, Cauchy--Schwarz gives
\[
 \frac94t^2\le t(27-3\alpha^2-4t),
 \qquad t\le\frac{108}{25}<5.
\]
Also $t\equiv\sum_i d_i=0\pmod2$. It remains to consider $t=0,2,4$.

If $t=4$, the $d_i$ are, in some order, $(1,1,-1,-1)$.
Let $P$ and $N$ be the sums of the $c_i$ at the positive and negative
positions. Then $P-N=-6$ and $P+N=3-\alpha$.
For $\alpha=0$ these equations contradict integrality.
For $\alpha=1$ they give $(P,N)=(-2,4)$, so
$\sum_i c_i^2\ge2+8=10$, whereas
\eqref{eq:cw65-reduced} requires $8$.
For $\alpha=-1$ they give $(P,N)=(-1,5)$, so
$\sum_i c_i^2\ge1+13=14$, again exceeding $8$.
Thus $t\ne4$.

If $t=2$, arrange $d=(1,-1,0,0)$. Then $c_1-c_2=-3$, so
\[
 (c_1,c_2)\in\{(-3,0),(-2,1),(-1,2),(0,3)\}.
\]
Their squared norms are $9,5,5,9$.
For $\alpha=\pm1$, the remaining squared norm $c_3^2+c_4^2$ must be
$7$ or $11$, neither a sum of two squares. For $\alpha=0$ it must be
$10$ or $14$. The latter is impossible. In the former case the first
pair is $(-3,0)$ or $(0,3)$, requiring $c_3+c_4=6$ or $0$.
But the solutions to $c_3^2+c_4^2=10$ have sums in $\{-4,-2,2,4\}$.
Thus $t\ne2$.

We conclude that $t=0$, so all $d_i=0$. If $\alpha=\pm1$, then
$\sum_i c_i^2=24$. Reduction modulo $4$ forces either zero or four odd
$c_i$. Four odd entries give squared sum $4\pmod8$, while zero odd
entries, with $|c_i|\le3$, give squared sum at most $16$.
Both are impossible. Hence $\alpha=0$ and $\sum_i c_i^2=27$.
There must be three odd entries, since one odd and three even entries
have squared sum at most $9+12=21$. Modulo $8$, the remaining even entry
must be zero. The three odd entries must then all have absolute value $3$.
Their sum is $3$, so two equal $3$ and one equals $-3$.
This proves the stated divisibility of every coefficient of $B$.
\end{proof}

\begin{proof}[Proof of Theorem~\ref{thm:cw195}]
Suppose $A\in\mathbb Z[C_{195}]$ represents such a matrix, with
$\epsilon(A)=9$. Write $C_{195}=C_{65}\times C_3$ and contract the
$C_3$ coordinate. The result is an integer contraction $B$ satisfying
Lemma~\ref{lem:icw65-divisible}; hence $B=3C$ with $C$ having coefficients
$0,\pm1$.

Let $\omega$ be a primitive cube root of unity and put $\lambda=1-\omega$,
so $\lambda\overline\lambda=3$. Apply the nonprincipal character
$Y\mapsto\omega$ of $C_3$ to $A$, obtaining
$E\in\mathbb Z[\omega][C_{65}]$ with $EE^{(-1)}=81$.
Each three-entry fiber has principal sum $0$ or $\pm3$.
A fiber of sum $\pm3$ is constant and has nonprincipal image zero.
A nonzero zero-sum fiber is a permutation of $(1,-1,0)$ and has image
$\lambda$ times an Eisenstein unit. Thus
\begin{equation}\label{eq:cw195-eisenstein}
 E=\lambda D,\qquad DD^{(-1)}=27,
 \qquad D_x\in\{0,\pm1,\pm\omega,\pm\omega^2\}.
\end{equation}
In particular, $D$ has precisely $27$ nonzero coefficients.

Apply the generalized multiplier theorem with coefficient conductor $3$,
group $C_{65}$, norm $27$, and multiplier integer $133$.
We have $(27,65)=1$, $(133,195)=1$, and
\[
 133\equiv1\pmod3,\qquad133\equiv3\pmod{65}.
\]
The automorphism $\sigma_{133}$ fixes every prime above $3$ in
$\mathbb Q(\zeta_{195})$: its restriction to $\mathbb Q(\zeta_{65})$
is the Frobenius automorphism at $3$, and adjoining $\omega$ is totally
ramified at these primes, so each has a unique extension.
The theorem therefore gives $D^{(133)}=\eta X^jD$.
Since $\sigma_{133}$ is the identity on $\mathbb Z[\omega]$, augmentation
gives $\epsilon(D)=\eta\epsilon(D)$. By
\eqref{eq:cw195-eisenstein}, $|\epsilon(D)|^2=27$, so $\eta=1$.
Finally $(133-1,65)=1$, allowing a group translation after which
$D^{(133)}=D$. Equivalently, $D$ is constant on multiplication-by-$3$
orbits of $C_{65}$.

These orbits have sizes $1,4$, four copies of $3$, and four copies of $12$.
The support size $27$, reduced modulo $3$, forces the size-$1$ and size-$4$
orbits both to be absent. The remaining support equation is
\[
 12k+3l=27,\qquad0\le k,l\le4,
\]
whose only solution is $(k,l)=(2,1)$. Consequently
\[
 \epsilon(D)=12u+12v+3w=3(4u+4v+w)
\]
for Eisenstein units $u,v,w$. Its norm must be $27$, so
$|4u+4v+w|=\sqrt3$. If $v=-u$, this absolute value is $1$.
Otherwise $u+v$ is a nonzero Eisenstein integer, and therefore
$|u+v|\ge1$. The reverse triangle inequality then gives
\[
 |4u+4v+w|\ge4|u+v|-|w|\ge3.
\]
Neither alternative equals $\sqrt3$. This contradiction proves the theorem.
\end{proof}

\paragraph{Reproducibility and limits of the transfer.}
The proof is entirely analytic. Independent standard-library calculations
enumerate $3388$ contraction candidates satisfying augmentation and energy,
recovering eight normalized contractions in the two previously recorded
equivalence classes. A separate exact calculation checks all $864$ normalized
Eisenstein orbit assignments; four survive the first correlation equation
and none survive the second. These checks are redundant certificates rather
than hypotheses of the proof.

The successful transfer is the retention of character information inside
fibers, followed by a bounded-coefficient and augmentation obstruction.
It is not a general solution of the circulant Hadamard problem. In fact, for
every odd $u$, its order-$4$ quotient at candidate order $4u^2$ has the
feasible vector $(u,u,u,-u)$. This vector has the correct parity, coefficient
bounds, and flat quotient Fourier norm, and each entry lifts separately to
a sum of $u^2$ signs. Thus a test using only the maximal $2$-power quotient
and those constraints cannot exclude any such candidate order.

\section{Verification and limits}

\subsection{Checks supplied with this note}
The general proofs above do not rely on finite search. The supplementary
programs use exact arithmetic to check the arithmetic reductions independently.
For the array result, two implementations verify 10,813 shell inequalities
through $k=1000$, including the boundary $d^2=k$, and check the constructions
against the original one- and two-coordinate marginal equations. A separate
small-parameter enumeration reproduces all minimum multiplicity vectors;
an optional CP-SAT run checks $3\le k\le16$.

For the circulant result, an exhaustive orbit calculation recovers eight
normalized contractions in two equivalence classes. An independent
Eisenstein calculation excludes the normalized unit-orbit assignments.
These numerical checks are unnecessary for the final proof: the contraction
lemma and the augmentation inequality already exclude every possible matrix.

\subsection{What transfers, and what does not}
For the array problem, positivity makes a populated orbit contribute at
least its full cardinality. The moment equations then force enough mass
outside it. There is no need for a $2$-adic obstruction. For the circulant
problem, a quotient alone discards information: evaluating a nonprincipal
character of its kernel retains enough information to force a unit-valued
group-ring equation. The resulting orbit sizes conflict with the permitted
augmentation norm.

This does not justify a general implication from local integral feasibility
to a solution with a prescribed finite alphabet or with nonnegative
coefficients. Graph-flow, marginal-table, tiling, and Littlewood-polynomial
targets were considered separately; no resolution of their broad open
conjectures is asserted here. In particular, neither result proves the
Hadamard conjecture, the circulant Hadamard conjecture, or Tutte's flow
conjectures. The evidence for transfer consists of the two precise theorems
proved above, with the hypotheses stated there.

\begin{thebibliography}{9}
\bibitem{Karunanayaka2026}
R. C. Karunanayaka,
\emph{Catalan Bounds for Symmetric Strength-Two Orthogonal Arrays},
Journal of Combinatorial Designs, first published 18 September 2026.
\href{https://doi.org/10.1002/jcd.70039}{doi:10.1002/jcd.70039}.
Conjecture 1 and Theorem 6.

\bibitem{ArasuGordonZhang2021}
K. T. Arasu, D. M. Gordon, and Y. Zhang,
\emph{New nonexistence results on circulant weighing matrices},
Cryptography and Communications \textbf{13} (2021), 775--789.
\href{https://arxiv.org/abs/1908.08447v3}{arXiv:1908.08447v3}.
Tables 9 and 10.

\bibitem{Tan2018}
M. M. Tan, \emph{Group Invariant Weighing Matrices},
Designs, Codes and Cryptography \textbf{86} (2018), 2677--2702.
\href{https://arxiv.org/abs/1610.01914v4}{arXiv:1610.01914v4}.
Theorem 3.4 and Corollary 3.5(b).

\bibitem{Tan2026}
M. M. Tan,
\emph{Character and Multiplier Obstructions for Circulant Weighing Matrices},
preprint, 1 September 2026.
\href{https://arxiv.org/abs/2608.18468v2}{arXiv:2608.18468v2}.

\end{thebibliography}
\end{document}
